% ============================================================================
% Tablet Deployment — report section + appendix
% ----------------------------------------------------------------------------
% Preamble requirement: \usepackage{listings}
% MagicMirror^2 uses \textsuperscript{2} (swap for $^2$ if your template prefers).
%
% This file contains TWO blocks:
%   1. MAIN TEXT  — a short \subsection for the body of the report.
%   2. APPENDIX   — the detailed procedure; place it after \appendix.
% The cross-reference resolves automatically once both are included.
% ============================================================================


% ============================== 1. MAIN TEXT ================================
\subsection{Tablet Deployment Script}
\label{sec:tablet-deployment}

The tablet --- which hosts both the system's MQTT broker and the
MagicMirror\textsuperscript{2} interface --- is provisioned by a single on-device
script (\texttt{deploy.sh}) that installs the broker, the
MagicMirror\textsuperscript{2} interface and the supporting services, and registers
them to start automatically whenever the tablet boots. A second, short step run
from the workstation then delivers the tablet's mutual-TLS certificates, which for
security reasons are never generated on the device itself. The full procedure and
its design rationale are described in Appendix~\ref{app:tablet-deploy}.


% ============================== 2. APPENDIX =================================
% Place this after \appendix in your report.
\section{Tablet Deployment Procedure}
\label{app:tablet-deploy}

Provisioning the tablet is more involved than a conventional Linux service
installation for two reasons. First, Android is a difficult host for unattended
background services: it has no service manager, applies aggressive power management
(Doze), and may terminate background processes under memory pressure. Second, the
system's security model forbids the project's certificate-authority (CA) private
key from ever residing on the tablet. The procedure below is shaped by these two
constraints.

\subsection{A two-stage procedure}

Provisioning is split into an on-device stage and a workstation stage. The
on-device stage is a single, idempotent script (\texttt{deploy.sh}) run once inside
Termux from a clone of the project repository; it installs and configures
everything that can be done locally on the tablet. The workstation stage then runs
from the developer laptop and pushes the tablet's mutual-TLS identity --- its broker
and mirror certificates together with the broker configuration --- onto the device
over SSH. This split follows directly from the security model: the CA private key
is the root of trust for the entire device fleet, so it must not be placed on an
always-on, physically accessible tablet. The tablet therefore receives only its own
leaf certificates and the public CA certificate, and the on-device script handles
no secret material.

\subsection{On-device provisioning (\texttt{deploy.sh})}

The script is idempotent (safe to re-run) and performs the following:
\begin{enumerate}
  \item Requests storage access and updates the Termux package index.
  \item Installs the required packages:
    \begin{itemize}
      \item \textbf{Git, Node.js} --- MagicMirror\textsuperscript{2} and its installer;
      \item \textbf{Mosquitto} --- the MQTT broker;
      \item \textbf{OpenSSH} --- remote access for the workstation stage;
      \item \textbf{Python} with the \texttt{zeroconf} and \texttt{paho-mqtt}
            libraries --- the mDNS responder and the presence-to-screen bridge;
      \item the \textbf{Termux:API} tools --- the wake-lock and the commands used
            to launch the kiosk browser.
    \end{itemize}
  \item Installs MagicMirror\textsuperscript{2} and the MQTT-client dependency of
        the custom bridge module. Because the script runs from a Git checkout,
        MagicMirror's installer --- which performs a clean-up step and fetches its
        bundled front-end assets --- completes correctly.
  \item Copies the runtime scripts into place and installs \texttt{tablet\_boot.sh}
        as the Termux:Boot hook.
  \item Starts the SSH daemon and prints the tablet's address together with the
        command to run on the workstation.
  \item Stages the kiosk-browser configuration and the companion Android
        application in the Downloads folder for installation.
\end{enumerate}
Until the certificates are delivered by the workstation stage, the mirror interface
loads but the broker remains down; this is intentional, as the broker cannot start
without its key material and the system does not fall back to an insecure mode.

\subsection{Self-healing autostart (\texttt{tablet\_boot.sh})}

Because Android offers no service manager, persistence is achieved with the
\emph{Termux:Boot} add-on, which runs a hook script after the device boots and is
first unlocked. On each boot the hook executes, in order: acquire a wake-lock (to
survive Doze); wait for Wi-Fi association; start the SSH daemon; start Mosquitto
(once its certificates are present); start the mDNS responder; start the
presence-to-screen bridge; and start MagicMirror\textsuperscript{2} with its kiosk
viewer. Each service is guarded so that re-running the hook never starts a
duplicate, and the mirror launcher restarts the interface with exponential back-off
on failure, so the tablet recovers from a reboot or power interruption without
intervention.

\subsection{Network-independent broker discovery}

The sensor modules must locate the broker without hard-coded addresses, as the home
network assigns IP addresses dynamically. The tablet advertises itself over
multicast DNS under the stable name \texttt{tablet.local}; each sensor module is
provisioned once with this name and resolves the current address at every
connection. The installation therefore tolerates address changes and network moves
without any re-provisioning of the sensor devices.

\subsection{Manual steps (Android restrictions)}

A few actions cannot be automated because of Android's security model and must be
performed once by the installer:
\begin{itemize}
  \item Disable battery optimisation for Termux (otherwise Doze suspends the services).
  \item Install the Termux:Boot add-on (from F-Droid).
  \item Unlock the device once after boot to arm the hook (it fires only after the
        first unlock).
  \item Import the kiosk-browser configuration and install the companion
        application when Android prompts.
\end{itemize}

\subsection{Summarised procedure}

\begin{lstlisting}[
  basicstyle=\ttfamily\small,
  breaklines=true,
  frame=single,
  captionpos=b,
  caption={Tablet deployment commands.},
  label={lst:tablet-deploy}]
On the tablet (once, in Termux):
  1.  pkg install -y git
  2.  git clone <repository-url>
  3.  cd <repository> && ./deploy.sh

On the workstation (once):
  4.  scripts/setup_tablet_ssh.sh        <tablet-ip>
  5.  scripts/deploy_home_kit.sh tablet  <home-id> <tablet-ip>

Manual Android steps (once): disable battery optimisation for Termux,
install Termux:Boot, unlock after boot, import the kiosk configuration.
\end{lstlisting}
